\documentclass[../main.tex]{subfiles}
\begin{document}

\chapter{分岐}
\lessoninfo{
  video={Lesson 4，動画 1--48},
  textbook={H\&P \cite{hennessy2017} §3.3，§3.9，§C.2},
  keywords={分岐予測，予測ミスペナルティ，不成立予測，分岐先バッファ，分岐履歴テーブル，2 ビット予測器，履歴ベース予測器，PShare，GShare，トーナメント予測器，階層型予測器，リターンアドレススタック，プリデコード},
}

第3章では，制御依存がどのパイプラインでもハザードになることを見た．
プロセッサは，分岐の行き先が分かるよりずっと前に分岐の後続命令をフェッチ
しなければならず，選択を誤るたびに作業がフラッシュされる．この章では，
フェッチステージがその選択をどのように行うかを調べる．最も単純な方針
---常に次の命令に進む---から出発して，実際のプロセッサで使われる構造を
順に導入する．すなわち，分岐先バッファ，飽和カウンタ，分岐履歴から
パターンを学習する予測器，複数の予測器の組合せ，そして関数からのリターン
専用の予測器である．その過程で，パイプラインが深いほど予測精度が重要に
なる理由も示す．

\section{パイプラインにおける分岐}
\label{sec:l04-branch-pipeline}

第3章の 5 段パイプラインで，条件分岐 \texttt{BEQ R1,R2,L} を考える．
\source{Lesson 4，動画 1--3；H\&P §C.2．}R1 と R2 が等しければ，実行は
ラベル~L から続く．このとき分岐は\term[せいりつぶんき]{成立分岐}[taken branch]%
であり，L のアドレスがその\term[ぶんきさきあどれす]{分岐先アドレス}[branch target address]%
である．

そうでなければ，実行はメモリ上の次の命令，すなわち PC+4 から続き，分岐は
\term[ふせいりつぶんき]{不成立分岐}[not-taken branch]となる．どちらの場合
であるかをパイプラインが知るのは，分岐が 3つのステージを通過した後である．
IF が分岐をフェッチし，ID が R1 と R2 を読み，EX がそれらを比較して分岐先
を計算する．分岐は EX の終わりに\emph{解決}される
（\cref{fig:l04-branch-pipeline}）．

フェッチステージはそれを待てない．分岐が ID にあるサイクル~2 にも，続く
サイクル~3 にも，フェッチステージは\emph{何らかの}命令を読み込まなければ
ならない．分岐が解決されるとき，その後ろにはすでに 2 命令がパイプライン
に入っている．それらが本当に分岐の後に続く命令であれば何も失われないが，
そうでなければフラッシュされ，2 サイクルが無駄になる．この間に何も
フェッチしなければ同じ 2 サイクルが無条件に無駄になるので，正しい可能性が
最も高い命令をフェッチすることが常に得策である．

\begin{figure}[htbp]
  \centering
  % A branch in the five-stage pipeline: the two instructions fetched while the
% branch is in ID and EX are predictions, checked when the branch resolves.
\begin{tikzpicture}[x=0.95cm,y=0.55cm, font=\footnotesize\sffamily,
  st/.style={draw, minimum width=0.9cm, minimum height=0.5cm, inner sep=0pt, fill=ln-box},
  pr/.style={st, fill=ln-accent!15}]
  \foreach \c in {1,...,6} { \node at (\c,0.8) {\c}; }
  \node[anchor=east] at (0.5,0.8) {\iflangja{サイクル}{cycle}};
  \node[anchor=east] at (0.5,0) {\texttt{BEQ R1,R2,L}};
  \foreach \s/\lab in {0/IF,1/ID,2/EX,3/MEM,4/WB} { \node[st] at (\s+1,0) {\lab}; }
  \node[anchor=east] at (0.5,-1) {\iflangja{予測した命令 1}{predicted instr.\ 1}};
  \node[pr] at (2,-1) {IF}; \node[pr] at (3,-1) {ID};
  \node[anchor=east] at (0.5,-2) {\iflangja{予測した命令 2}{predicted instr.\ 2}};
  \node[pr] at (3,-2) {IF};
  \draw[ln-caution, thick, dashed] (3.5,1.15) -- (3.5,-2.6)
    node[below, font=\scriptsize] {\iflangja{分岐解決（EX の終わり）}{branch resolved (end of EX)}};
  \node[anchor=west, align=left, font=\scriptsize, text=ln-muted] at (3.75,-1.5)
    {\iflangja{予測が正しい: 両命令とも続行\\予測ミス: 両命令をフラッシュ}{correct prediction: both continue\\misprediction: both are flushed}};
\end{tikzpicture}

  \caption{5 段パイプラインにおける分岐．分岐が解決される前にフェッチ
    される 2 命令は予測であり，予測が誤っていればフラッシュされる．}
  \label{fig:l04-branch-pipeline}
\end{figure}

\begin{definition}[分岐予測]
  \label{def:l04-branch-prediction}
  現在の命令がデコードも実行もされないうちに，次の命令をフェッチする
  アドレスをフェッチステージで選ぶことを\term[ぶんきよそく]{分岐予測}[branch prediction]%
  という．
\end{definition}

\begin{definition}[予測ミス]
  \term[ぶんきよそくみす]{分岐予測ミス}[branch misprediction]は，解決
  された分岐が予測とは別の場所へ進むと判明したときに生じる．誤った
  アドレスからフェッチされた命令はフラッシュされ，正しいアドレスから
  フェッチがやり直される．
\end{definition}

\subsection{予測器が知っていること}

フェッチの時点で得られる情報は，フェッチ中の命令の PC だけである．命令
そのものはまだデコードすらされていない．このアドレスだけから，予測器は
3つの問いに答えなければならない．
\begin{enumerate}
  \item その命令はそもそも分岐か．
  \item 分岐であれば，成立するか．
  \item 成立するならば，分岐先アドレスはどこか．
\end{enumerate}
答えは次の PC という 1つの数である．分岐でない命令と不成立分岐では
PC+4，成立分岐では分岐先アドレスとなる．\margin{無条件ジャンプ，コール，
リターンは常に成立する分岐であり，これらに関係するのは問い 1 と~3 だけ
である．}

\section{予測精度と性能}
\label{sec:l04-accuracy}

予測器があれば，制御ハザードで失われるサイクルは，予測ミスした経路で
費やされたサイクルだけになる．\source{Lesson 4，動画 4--5；H\&P §C.2．}%
\cref{eq:pipeline-cpi} を分岐に特化すると次式を得る．
\begin{equation}
  \text{CPI} \;=\; 1 \;+\;
  \frac{\text{予測ミス数}}{\text{命令}} \times
  \frac{\text{ペナルティ}}{\text{予測ミス}}
  \label{eq:l04-cpi}
\end{equation}

\begin{definition}[精度とペナルティ]
  \term[ぶんきよそくせいど]{分岐予測精度}[branch prediction accuracy]%
  とは，結果が正しく予測される分岐の割合である．したがって
  \begin{equation*}
    \frac{\text{予測ミス数}}{\text{命令}} =
    \frac{\text{分岐数}}{\text{命令}} \times
    (1 - \text{精度})
  \end{equation*}
  となる．ここで，分岐でない命令は次の PC を PC+4 として常に正しく予測
  されると仮定している．\term[よそくみすぺなるてぃ]{予測ミスペナルティ}[misprediction penalty]%
  とは，予測ミス 1 回当たりに失われるサイクル数である．1 ステージに 1 命令
  を保持するパイプラインでは，フラッシュされる命令数に等しい．ステージ
  $r$ で解決される分岐のペナルティは $r-1$ サイクルである．
\end{definition}

\cref{eq:l04-cpi} の 2つの因子は由来が異なる．第 1 の因子は予測器（および
プログラムが実行する分岐の数）で決まり，第 2 の因子はパイプライン，すなわち
分岐がどれだけ遅く解決されるかで決まる．よりよい予測器は第 1 の因子を
小さくし，より深いパイプラインは第 2 の因子を大きくする．そのため同じ
予測器でも，深いパイプラインほどコストが大きくなる．完全な予測器であれば，
深さによらず理想の CPI~1 に到達する．

\begin{example}
  \label{ex:l04-cpi}
  あるプログラムでは命令の \SI{15}{\percent} が分岐であり，予測器の精度は
  \SI{92}{\percent} である．命令当たりの予測ミス数は
  $0.15 \times 0.08 = 0.012$ である．分岐が EX（ステージ~3，ペナルティ~2）
  で解決されるなら $\text{CPI} = 1 + 0.012 \times 2 = 1.024$ となる．
  分岐をステージ~11（ペナルティ~10）で解決するパイプラインでは，同じ
  予測器で $\text{CPI} = 1 + 0.012 \times 10 = 1.12$ となる．
\end{example}

\section{不成立予測}
\label{sec:l04-not-taken}

予測に代わる方法は，予測を拒むことである．フェッチステージは次の PC が
分かるまで単に待つ．\source{Lesson 4，動画 6--8；H\&P §C.2．}命令が分岐
かどうかさえ ID の後にならないと分からないので，この待ちはすべての命令に
影響する．分岐でない命令の次のフェッチはその命令自身のフェッチの 2 サイクル
後になり，分岐の次のフェッチは，成立・不成立にかかわらず 3 サイクル後に
なる．

\begin{table}[b]
  \centering\small
  \caption{5 段パイプラインにおいて，ある命令のフェッチから正しい次の命令
    がフェッチされるまでのサイクル数．}
  \label{tab:l04-not-taken}
  \begin{tabular}{@{}lcc@{}}
    \toprule
    命令          & 予測を拒む & 不成立予測 \\
    \midrule
    成立分岐      & 3 & 3 \\
    不成立分岐    & 3 & 1 \\
    分岐以外      & 2 & 1 \\
    \bottomrule
  \end{tabular}
\end{table}

最も単純な実用的予測器は\term[ふせいりつよそく]{不成立予測}[predict not taken]，
すなわち常に PC+4 からフェッチすることである．フェッチステージはどのみち
PC をインクリメントするので，これにはメモリも，フェッチステージがすでに
持つもの以外のハードウェアも必要ない．分岐でない命令と不成立分岐に対しては
常に正しく，誤るのは成立分岐だけであり，そのコストは予測を拒む場合と同じく
3 サイクルである．\cref{tab:l04-not-taken} は 2つの方針を比較している．不成立予測が
劣ることは決してない．

しかし，その精度は低い．大半の分岐は成立する---ループの分岐はたいてい
後方へジャンプし，無条件ジャンプ，コール，リターンは常に成立する
---ので，不成立予測は分岐の過半数を予測ミスする．第3章で用いた典型的な
統計では，全命令の \SI{12}{\percent} で予測ミスが起こり，
\cref{eq:branch-cpi} の CPI となる．

\begin{note}[処理中の複数の予測]
  深いパイプラインでは，IF と解決ステージの間に複数の分岐が同時に存在
  しうる．それぞれの分岐は，フェッチ時に行われた予測を伴っている．ある
  分岐が予測ミスと判明すると，その後にフェッチされたものは，後続の分岐と
  その予測も含めてすべてフラッシュされる．したがって，誤った経路で
  フェッチされた分岐がそれ自身のフラッシュを引き起こすことはない．
  サイクルを失わせるのは，正しい経路上にある分岐の予測ミスだけである．
\end{note}

\section{よりよい予測の必要性}
\label{sec:l04-better}

よりよい予測器にどれだけの価値があるかは，パイプラインに依存する．
\source{Lesson 4，動画 9--11；H\&P §3.3．}1 ステージ当たり $w$ 命令を扱う
プロセッサは，同じ深さで幅が 1 命令のパイプラインと比べて，予測ミス 1 回
当たり $w$ 倍の命令をフラッシュする．したがって，完全な予測に対する相対
実行時間は
\begin{equation}
  \frac{T}{T_{\text{perfect}}} \;=\; 1 +
  \frac{\text{予測ミス数}}{\text{命令}} \times
  (\text{予測ミス当たりのフラッシュ命令数})
  \label{eq:l04-relative}
\end{equation}
であり，$w=1$ のときはちょうど\cref{eq:l04-cpi} の CPI になる．\margin{実際
の値：Skylake 級のコアは予測ミス 1 回につき約 17 サイクルを失い，224 エントリ
のリオーダバッファを埋め直す（Intel 最適化マニュアル）．}

\begin{example}
  \label{ex:l04-better}
  命令の \SI{15}{\percent} が分岐で，そのうち \SI{70}{\percent} が成立
  するとする．このとき不成立予測の精度は \SI{30}{\percent} である．
  \cref{tab:l04-better} は，3 種類の予測器と 3 種類のパイプラインについて
  \cref{eq:l04-relative} を評価したものである．パイプラインは，5 段
  パイプライン（フラッシュ 2 命令），分岐をステージ~13 で解決する
  パイプライン（フラッシュ 12 命令），および同じ深さで 1 ステージ当たり
  3 命令を扱うパイプライン（フラッシュ 36 命令）である．
  \caution{予測ミスは分岐単位ではなく\emph{命令}単位で数える．$1-\text{精度}$
  に，命令のうち分岐が占める割合を掛ける．}
\end{example}

\begin{table}[htbp]
  \centering\small
  \caption{\cref{ex:l04-better} の分岐統計における，完全な予測に対する
    相対実行時間．}
  \label{tab:l04-better}
  \begin{tabular}{@{}lcccc@{}}
    \toprule
    予測器 & 予測ミス/命令 & ステージ 3 & ステージ 13 & ステージ 13，幅 3 \\
    \midrule
    不成立（\SI{30}{\percent}） & 0.105  & 1.21  & 2.26  & 4.78 \\
    精度 \SI{90}{\percent}      & 0.015  & 1.03  & 1.18  & 1.54 \\
    精度 \SI{99}{\percent}      & 0.0015 & 1.003 & 1.018 & 1.054 \\
    \bottomrule
  \end{tabular}
\end{table}

ここから 3つのことが分かる．\tip{（命令当たりの予測ミス数）$\times$
（フラッシュ命令数）がそのまま性能低下率になる．$0.015 \times 12 = 0.18$，
すなわち中段の \SI{18}{\percent} である．}浅いパイプラインでは，不成立予測でさえコストは
\SI{21}{\percent} にすぎず，精度を \SI{90}{\percent} から
\SI{99}{\percent} に改善しても利得は \SI{3}{\percent} に満たない．深い
パイプラインでは，不成立予測は実行時間を 2 倍以上にし，同じ改善は
$1.18/1.018 \approx 1.16$ の高速化に相当する．深く幅も広いパイプライン
では $1.54/1.054 \approx 1.46$ に相当する．現代のプロセッサは深く，かつ
幅も広い．だからこそ分岐予測に多大な投資をしているのである．

\begin{keypoint}
  重要なのは精度ではなく\emph{予測ミス}率である．精度 \SI{90}{\percent}
  から \SI{99}{\percent} への向上は小さな一歩に見えるが，予測ミスの
  \SI{90}{\percent} を取り除く．そしてパイプラインが深く幅広いほど，残った
  予測ミス 1 回当たりのコストは大きくなる．
\end{keypoint}

\section{分岐先バッファ}
\label{sec:l04-btb}

不成立予測は，フェッチ中の命令に関する情報を何も使わない．よりよい予測器
はその\emph{履歴}，すなわちこの PC にある命令が過去に実行されたときに何を
したかを用いる．\source{Lesson 4，動画 12--15；H\&P §3.9；\cite{lee1984}．}%
次の PC は，PC とこの履歴の関数として予測される．あるアドレスに格納された
命令は変化しないので，それが分岐かどうか，直接分岐がどこへジャンプするか
は毎回同じである．また，ほとんどの分岐はたいてい同じ方向に進む．

\begin{definition}[分岐先バッファ]
  \term[ぶんきさきばっふぁ]{分岐先バッファ}[branch target buffer]%
  \index{BTB|see{分岐先バッファ}}（BTB）とは，フェッチ中の命令の PC で
  インデックスされ，その命令の予測された次の PC を格納するテーブルである．
\end{definition}

IF では現在の PC で BTB を読み，その出力を次の PC として用いる．分岐が
解決されると，実際の次の PC を予測と比較する．一致しなければ，誤って
フェッチされた命令をフラッシュし，実際の分岐先を BTB のエントリに書き込む．
こうして，同じ命令の次の実行は新しい分岐先で予測される．

\subsection{現実的な BTB}

BTB は，あり得るすべての命令アドレスに 1つずつエントリを持つことはできない．
BTB はフェッチそのものに加えて同じフェッチサイクル内で読まなければならない
ので，小さなテーブル---数百から数千エントリ---に限られる．\margin{実際の
容量：Intel Skylake の BTB は約 4096 エントリ，AMD の Zen~2 は 512 エントリと
7168 エントリの 2 段構成である（各社の最適化ガイド）．}そこで PC は
単純なハッシュ，すなわちその下位の数ビットによってエントリに対応付けられる．
ループの命令のように時間的に近接して実行される命令は，メモリ上でも近接して
いる．それらは上位ビットを共有するが下位ビットが異なるので，別々のエントリ
に分散される．2つの命令が衝突するのはコード上で遠く離れている場合だけで
あり，そのような命令が同時に活発に実行される可能性は低い．

ワード境界に整列した 4 バイト固定長命令では，すべての PC の下位 2 ビット
はゼロであり情報を持たない．そのため $2^k$ エントリの BTB はビット $k+1$
から~2 でインデックスされる（\cref{fig:l04-btb}）．\tip{エントリ番号 $=
(\text{PC} \gg 2) \mathbin{\&} (\text{エントリ数} - 1)$．常にゼロの 2 ビットを
捨て，下位 $k$ ビットを残す．}可変長命令（x86）では，
インデックスはビット~0 から始まる．

\begin{figure}[htbp]
  \centering
  % A realistic BTB: low-order PC bits (above the alignment bits) index a
% small table of predicted targets.
\begin{tikzpicture}[x=1cm,y=1cm, font=\footnotesize\sffamily, >={Stealth[length=1.8mm]}]
  % PC fields
  \node[anchor=east] at (-0.1,0.3) {PC};
  \draw[fill=white] (0,0) rectangle (3.0,0.6);   \node at (1.5,0.3) {31 \dots\ 10};
  \draw[fill=ln-accent!15] (3.0,0) rectangle (5.6,0.6); \node at (4.3,0.3) {9 \dots\ 2};
  \draw[fill=ln-box] (5.6,0) rectangle (6.4,0.6); \node at (6.0,0.3) {1\,0};
  \node[anchor=west, font=\scriptsize, text=ln-muted] at (6.5,0.3) {\iflangja{常に 00}{always 00}};
  % index
  \draw[->] (4.3,0) -- (4.3,-0.8) node[midway, right, font=\scriptsize] {\iflangja{インデックス（8 ビット）}{index (8 bits)}};
  % table: 6 rows of 0.45
  \fill[ln-accent!15] (3.0,-2.15) rectangle (5.6,-2.6);
  \foreach \i in {0,...,5} { \draw (3.0,-0.8-0.45*\i) rectangle (5.6,-1.25-0.45*\i); }
  \draw (3.7,-0.8) -- (3.7,-3.5);
  \foreach \i/\lab in {0/0,1/1,3/170,5/255} {
    \node[font=\scriptsize, text=ln-muted] at (3.35,-1.025-0.45*\i) {\lab}; }
  \foreach \i in {2,4} { \node[font=\scriptsize, text=ln-muted] at (3.35,-1.0-0.45*\i) {$\vdots$}; }
  \node at (4.65,-2.375) {\iflangja{分岐先}{target}};
  \node[font=\scriptsize, text=ln-muted, anchor=north] at (4.3,-3.55) {BTB};
  % lookup and update
  \draw[->, thick] (5.6,-2.375) -- (7.0,-2.375) node[right] {\iflangja{予測した次の PC}{predicted next PC}};
  \draw[->, dashed, ln-caution] (0.6,-2.375) -- (3.0,-2.375)
    node[pos=0.0, above right, font=\scriptsize, align=left] {\iflangja{分岐解決時に\\実際の分岐先を書く}{on resolution:\\write actual target}};
\end{tikzpicture}

  \caption{PC のビット 9--2 でインデックスされる 256 エントリの BTB．
    エントリは IF で読まれ，分岐が解決されると実際の分岐先で書き込まれる．}
  \label{fig:l04-btb}
\end{figure}

\begin{example}
  \label{ex:l04-btb-index}
  256 エントリの BTB では，\texttt{0x004012A8} にある命令は PC のビット
  9--2 を用いる．$(\texttt{0x12A8} \gg 2) \mathbin{\&} \texttt{0xFF} =
  \texttt{0x4AA} \mathbin{\&} \texttt{0xFF} = \texttt{0xAA}$，すなわち
  エントリ~170 である．次の命令（\texttt{0x004012AC}）はエントリ~171 を
  用い，256 命令先の \texttt{0x004016A8} にある命令はエントリ~170 を共有
  する．
\end{example}

\section{方向予測}
\label{sec:l04-direction}

BTB だけでも，方向の問いには暗黙のうちに答えている．PC+4 を保持する
エントリは「不成立」を意味する．\source{Lesson 4，動画 16--21；H\&P §C.2．}%
これでは，分岐でない命令と不成立分岐のすべてに分岐先アドレス 1つ分の
記憶を費やすことになる．予測を分割する方が安価である．
\term[ほうこうよそくき]{方向予測器}[direction predictor]が，その命令が成立
分岐であるかどうかを判断し，成立分岐と判断したときにだけ BTB を参照する．

\begin{definition}[分岐履歴テーブル]
  \term[ぶんきりれきてーぶる]{分岐履歴テーブル}[branch history table]%
  \index{BHT|see{分岐履歴テーブル}}（BHT）とは，PC の下位ビットで
  インデックスされるテーブルの形をとる方向予測器である．最も単純な形では，
  各エントリは分岐の前回の実行結果を保持する 1 ビットであり，0 なら次の PC
  として PC+4 を，1 なら BTB の分岐先を選ぶ．これが
  \term[1びっとよそくき]{1 ビット予測器}[one-bit predictor]\index{1BP|see{1 ビット予測器}}%
  （1BP）である．
\end{definition}

分岐が解決されると，その BHT エントリには実際の結果が設定され，分岐が
成立した場合には分岐先が BTB に書き込まれる（\cref{fig:l04-bht-btb}）．
この構成には 3つの利点がある．分岐でない命令の BHT エントリは書き込まれる
ことがなく 0 のままなので，PC+4 を正しく予測する．\caution{どちらのテーブル
にもタグがない．成立分岐が同じエントリに割り当てられると，分岐でない命令も
成立と予測され，誤りはデコード時にしか分からない．}BTB に必要なのは成立
分岐のためのエントリだけであり，成立分岐は全命令よりはるかに少ない．
そして BHT のエントリは 1 ビットなので，同じコストで BHT は BTB よりはるか
に多くのエントリを持つことができ，異なる分岐がエントリを共有することは
まれになる．

\begin{figure}[htbp]
  \centering
  % Fetch-stage predictor: the BHT predicts the direction, the BTB the target.
\begin{tikzpicture}[x=1cm,y=1cm, font=\footnotesize\sffamily, >={Stealth[length=1.8mm]},
  box/.style={draw, fill=ln-box, minimum height=0.6cm, inner sep=2pt}]
  \node[box, minimum width=0.9cm] (pc) at (0,0) {PC};
  % index bus
  \draw[->] (pc.east) -- (6.0,0) -- (6.0,-0.5);
  \draw[->] (1.2,0) -- (1.2,-0.9);
  \draw[->] (3.4,0) -- (3.4,-0.5);
  \fill (1.2,0) circle (1.2pt); \fill (3.4,0) circle (1.2pt);
  % +4
  \node[box, minimum width=0.8cm] (inc) at (1.2,-1.2) {$+4$};
  % BTB
  \draw[fill=ln-box] (2.5,-0.5) rectangle (4.3,-1.9);
  \foreach \y in {-0.78,-1.06,-1.34,-1.62} { \draw[ln-rule] (2.5,\y) -- (4.3,\y); }
  \node[anchor=east, font=\footnotesize\bfseries] at (2.45,-0.75) {BTB};
  % BHT
  \draw[fill=ln-box] (5.7,-0.5) rectangle (6.3,-1.9);
  \foreach \y in {-0.78,-1.06,-1.34,-1.62} { \draw[ln-rule] (5.7,\y) -- (6.3,\y); }
  \node[anchor=west, font=\footnotesize\bfseries] at (6.35,-0.75) {BHT};
  \node[anchor=west, font=\scriptsize, text=ln-muted] at (6.35,-1.15) {\iflangja{1 ビット/エントリ}{1 bit/entry}};
  % update
  \node[font=\scriptsize, text=ln-caution, align=center] (upd) at (5.0,-2.3) {\iflangja{EX で\\更新}{update\\from EX}};
  \draw[->, dashed, ln-caution] (upd.north) -- (4.3,-1.2);
  \draw[->, dashed, ln-caution] (upd.north) -- (5.7,-1.2);
  % mux
  \draw[fill=white] (0.7,-2.6) -- (3.9,-2.6) -- (3.3,-3.2) -- (1.3,-3.2) -- cycle;
  \draw[->] (inc.south) -- (1.2,-2.6);
  \draw[->] (3.4,-1.9) -- (3.4,-2.6);
  \node[font=\scriptsize] at (1.3,-2.8) {0};
  \node[font=\scriptsize] at (3.3,-2.8) {1};
  \draw[->] (6.0,-1.9) -- (6.0,-2.9) -- (3.6,-2.9)
    node[pos=0.72, below, font=\scriptsize] {\iflangja{成立?}{taken?}};
  % next PC feedback
  \draw[->] (2.3,-3.2) -- (2.3,-3.7) -- (-0.9,-3.7) -- (-0.9,0) -- (pc.west);
  \node[below, font=\scriptsize] at (0.7,-3.7) {\iflangja{次の PC}{next PC}};
\end{tikzpicture}

  \caption{BHT と BTB による次の PC の予測．BHT のビットが PC+4 と BTB の
    分岐先のどちらかを選ぶ．分岐が解決されると両方のテーブルが更新される．}
  \label{fig:l04-bht-btb}
\end{figure}

\begin{example}
  \label{ex:l04-1bp-trace}
  ある内側ループは，入るたびに配列の 3 要素を処理する．
  \begin{lstlisting}[language={[x86masm]Assembler}, numbers=none]
0x0400  Inner: LW   R2, 0(R3)     ; 要素をロード
0x0404         ADD  R5, R5, R2    ; 合計に加算
0x0408         ADDI R3, R3, 4     ; 次の要素
0x040C         ADDI R1, R1, -1    ; カウントを減らす
0x0410         BNE  R1, R0, Inner ; 1 回の実行で T, T, N
\end{lstlisting}
  64 エントリの BHT（インデックスはビット 7--2）では，5つの命令は BHT
  エントリ 0 から~4 を用い，\texttt{BNE} はエントリ~4 を用いる．8 エントリ
  の BTB（ビット 4--2）では，\texttt{BNE} は BTB エントリ~4 を用い，これを
  8 命令前の \texttt{0x03F0} にある命令と共有する．その命令が分岐でなければ，
  その BHT エントリ（エントリ~60）は 0 のままなので，BTB エントリを使うのは
  この分岐だけである．\cref{tab:l04-1bp-trace} は，すべてのテーブルが
  クリアされた状態から始めて，ループ 2 回の実行にわたり \texttt{BNE} を
  追跡したものである．
\end{example}

\begin{table}[htbp]
  \centering\small
  \caption{\cref{ex:l04-1bp-trace} の \texttt{BNE}（BHT エントリ~4）の，
    ループ 2 回の実行にわたる 1 ビット予測．}
  \label{tab:l04-1bp-trace}
  \begin{tabular}{@{}lcccccc@{}}
    \toprule
    実行              & 1 & 2 & 3 & 4 & 5 & 6 \\
    \midrule
    実行前の BHT[4]   & 0 & 1 & 1 & 0 & 1 & 1 \\
    予測した次の PC   & \texttt{0x414} & \texttt{0x400} & \texttt{0x400} & \texttt{0x414} & \texttt{0x400} & \texttt{0x400} \\
    結果              & T & T & N & T & T & N \\
    正しいか          & いいえ & はい & いいえ & いいえ & はい & いいえ \\
    \bottomrule
  \end{tabular}
\end{table}

最初の実行では，BTB エントリ~4 に \texttt{0x400} も書き込まれる．それ以降，
分岐先は常に正しいが，方向は 6 回の実行のうち 4 回で誤っている．

\section{2 ビット予測器}
\label{sec:l04-2bp}

1 ビット予測器は常に同じ方向に進む分岐にはよく機能するが，2つの弱点がある．
\source{Lesson 4，動画 22--26；H\&P §C.2；\cite{smith1981}．}
\begin{itemize}
  \item ほとんど常に一方向に進む分岐は，例外 1 回につき\emph{2 回}の予測
    ミスを被る．1 回は例外そのもので，もう 1 回はその次の通常の結果である．
    例外によってビットが反転してしまうからである．
  \item 方向が頻繁に変わる分岐はうまく予測できない．交互の分岐
    （T,~N,~T,~N,~\ldots）は毎回予測ミスし，\cref{tab:l04-1bp-trace} の
    短いループは 3 回に 2 回予測ミスする．
\end{itemize}

\begin{definition}[2 ビット予測器]
  \term[2びっとよそくき]{2 ビット予測器}[two-bit predictor]%
  \index{2BP|see{2 ビット予測器}}\index{2BC|see{2 ビット予測器}}（2BP）は，
  2 ビットカウンタ（2BC）とも呼ばれ，各 BHT エントリに，分岐が成立すると
  増加し，不成立だと減少する 2 ビットカウンタを格納する．
\end{definition}

このカウンタは\term[ほうわかうんた]{飽和カウンタ}[saturating counter]で
ある．すなわち，11 で増加させても 11 に留まり，00 で減少させても 00 に
留まる．上位ビットは\emph{予測ビット}であり，1 なら成立を予測する．下位
ビットは\emph{ヒステリシスビット}である．下位ビットは予測ビットと合わせて，
カウンタが自らの予測をどれだけ強く信じているかを記録する．2 ビットが
等しいときカウンタは強い状態に，異なるときは弱い状態にあるので，4つの
状態は強い不成立（00），弱い不成立（01），弱い成立（10），強い成立（11）
である（\cref{fig:l04-2bc}）．\tip{2 ビットを 0--3 の数とみなすとよい．T で
1 増え，N で 1 減り，両端で飽和する．値が 2 以上なら成立と予測する．}強い状態
から予測を変えるには，予測に反する結果が 2 回続く必要があるので，単発の例外に
よる予測ミスは 1 回で済む．

\begin{figure}[htbp]
  \centering
  % Two-bit saturating counter: T increments, N decrements.
\begin{tikzpicture}[font=\footnotesize\sffamily, >={Stealth[length=1.8mm]},
  s/.style={circle, draw, minimum size=0.95cm, inner sep=1pt},
  nt/.style={s, fill=ln-box}, tk/.style={s, fill=ln-accent!15}]
  \node[nt] (s0) at (0,0)   {00};
  \node[nt] (s1) at (2.5,0) {01};
  \node[tk] (s2) at (5.0,0) {10};
  \node[tk] (s3) at (7.5,0) {11};
  \draw[->] (s0) to[bend left=30] node[above] {T} (s1);
  \draw[->] (s1) to[bend left=30] node[above] {T} (s2);
  \draw[->] (s2) to[bend left=30] node[above] {T} (s3);
  \draw[->] (s3) to[bend left=30] node[below] {N} (s2);
  \draw[->] (s2) to[bend left=30] node[below] {N} (s1);
  \draw[->] (s1) to[bend left=30] node[below] {N} (s0);
  \draw[->] (s0) to[out=150,in=210,looseness=6] node[left] {N} (s0);
  \draw[->] (s3) to[out=30,in=-30,looseness=6] node[right] {T} (s3);
  \foreach \n/\lab in {s0/{\iflangja{強い不成立}{strong N}},s1/{\iflangja{弱い不成立}{weak N}},
                       s2/{\iflangja{弱い成立}{weak T}},s3/{\iflangja{強い成立}{strong T}}} {
    \node[font=\scriptsize, text=ln-muted, below=5mm of \n] {\lab};
  }
  \draw[ln-rule] (-0.6,-1.35) -- (3.1,-1.35);
  \draw[ln-rule] (4.4,-1.35) -- (8.1,-1.35);
  \node[font=\scriptsize, below] at (1.25,-1.35) {\iflangja{予測: 不成立}{predict not taken}};
  \node[font=\scriptsize, below] at (6.25,-1.35) {\iflangja{予測: 成立}{predict taken}};
\end{tikzpicture}

  \caption{2 ビット飽和カウンタの状態遷移図．T（成立）で右へ，N（不成立）
    で左へ移る．上位ビットが予測である．}
  \label{fig:l04-2bc}
\end{figure}

\begin{table}[htbp]
  \centering\small
  \caption{\cref{ex:l04-1bp-trace} のループ分岐に対する，00 から始まる
    2 ビットカウンタによる予測（ループ 3 回の実行）．予測ミスを $\times$
    で示す．}
  \label{tab:l04-2bc-trace}
  \setlength{\tabcolsep}{4.5pt}
  \begin{tabular}{@{}lccccccccc@{}}
    \toprule
    実行              & 1 & 2 & 3 & 4 & 5 & 6 & 7 & 8 & 9 \\
    \midrule
    結果              & T & T & N & T & T & N & T & T & N \\
    実行前のカウンタ  & 00 & 01 & 10 & 01 & 10 & 11 & 10 & 11 & 11 \\
    予測              & N & N & T & N & T & T & T & T & T \\
    2BP               & $\times$ & $\times$ & $\times$ & $\times$ & & $\times$ & & & $\times$ \\
    1BP（比較用）     & $\times$ & & $\times$ & $\times$ & & $\times$ & $\times$ & & $\times$ \\
    \bottomrule
  \end{tabular}
\end{table}

\cref{tab:l04-2bc-trace} は短いループに対する効果を示している．カウンタが
いったん強い成立（11，実行~6 の前）に達すると，1 回のループ脱出では
カウンタは弱い成立（10）に下がるだけで，次の成立の結果で元に戻る．その
ため予測ミスするのは脱出だけであり，予測ミスはループ 1 回の実行につき
2 回ではなく 1 回になる．\caution{\emph{弱い}状態からは，予測に反する結果
1 回で予測が変わる．2 ビットカウンタが予測を変えるのに，常に 2 回の予測ミスが
必要なわけではない．}

\subsection{初期化とビット数の増加}

強い状態に初期化されたカウンタは，分岐が実際には逆方向に進む場合に 2 回の
予測ミスを招く．弱い状態に初期化されたカウンタなら 1 回で済む．しかし弱い
状態は，カウンタが最も脆弱な状態でもある．交互の分岐 T,~N,~T,~N,~\ldots{}
は，カウンタが 01 から始まると 01 と 10 の間を行き来して毎回予測ミスするが，
00 から始まれば予測ミスは 1 回おきで済む．どちらの選択もすべての分岐に
とって最善ではなく，どの初期状態にも，それが不利になる結果の系列が存在
する．そこでカウンタは，単に最も単純な状態である 00 にリセットされる．
\margin{どの予測器も万能ではない．どんな予測器に対しても，常にその時点の
予測に反する結果の系列が存在する．}

3 ビット以上のカウンタはヒステリシスを増やす．すなわち，予測を変えるのに
より多くの例外が連続する必要がある．これが役立つのは例外がまとまって
生じる分岐だけであり，振る舞いを本当に変えた分岐への追従を遅くし，ビット
数も増える．1 ビットから 2 ビットへの変更は大きな改善だが，それ以上ビット
を増やしても得るものは少ない．2 ビットカウンタが標準的な構成要素である
のはこのためである．

\section{履歴ベース予測器}
\label{sec:l04-history}

分岐の中には，カウンタでは捉えられない規則的なパターンに従うものがある．
カウンタはどちらの方向が多いかを追跡するだけだからである．
\source{Lesson 4，動画 27--33；H\&P §3.3；\cite{yeh1991}．}ループ
インデックスが偶数かどうかを判定する分岐は T,~N,~T,~N と交互に変わり，
\cref{ex:l04-1bp-trace} のループ分岐は T,~T,~N を繰り返す．どちらも直近の
結果が分かれば完全に予測できる．

\begin{definition}[履歴ベース予測器]
  \term[りれきべーすよそくき]{履歴ベース予測器}[history-based predictor]%
  は，分岐の直近の結果（その履歴）を記録し，それを用いて複数のカウンタ
  から 1つを選ぶ．各カウンタは，ある特定の履歴の後に分岐が何をするかを
  学習する．
\end{definition}

\subsection{1 ビット履歴}

最も単純な構成では，各 BHT エントリに 1 ビットの履歴（分岐の前回の結果）と
2つの 2 ビットカウンタを格納する．一方は前回の結果が N のときに，
他方は T のときに使われる．予測では，履歴ビットがカウンタを選び，その
カウンタの上位ビットが予測となる．更新では，選ばれたカウンタを実際の
結果に従って増加または減少させ，その結果を新しい履歴ビットとする．
エントリには $1 + 2\times 2 = 5$ ビットが必要である．

交互に変わる分岐では，N の後に選ばれるカウンタは T だけを，T の後に
選ばれるカウンタは N だけを見る．それぞれの学習が済めば，分岐は毎回正しく
予測される．パターン T,~T,~N にはこれでは足りない．T の後で分岐は成立する
こともしないこともあるので，T の後に選ばれるカウンタは予測ミスを繰り返す．

\subsection{2 ビット履歴と \texorpdfstring{$N$}{N} ビット履歴}

履歴 2 ビットとカウンタ 4つ（エントリ当たり 10 ビット，
\cref{fig:l04-history}）では，可能な 4つの履歴 NN，NT，TN，TT のそれぞれが
専用のカウンタを持つ．T,~T,~N では，履歴 TT の後は常に N であり，TN と
NT の後は常に T である．\cref{tab:l04-history-trace} はループ分岐を追跡
したものである．カウンタの学習が済むと，すべての実行が正しく予測される．
NN のカウンタは二度と使われない．

\begin{figure}[htbp]
  \centering
  % Entry of a two-bit history predictor: the history selects one of four
% two-bit counters.
\begin{tikzpicture}[x=1cm,y=1cm, font=\footnotesize\sffamily, >={Stealth[length=1.8mm]}]
  \node[draw, fill=ln-box, minimum width=0.9cm, minimum height=0.6cm] (pc) at (-1.7,0) {PC};
  \draw[->] (pc.east) -- (0,0) node[midway, above, font=\scriptsize] {\iflangja{索引}{index}};
  \draw[fill=ln-accent!15] (0,-0.3) rectangle (1.4,0.3); \node at (0.7,0) {$h_1h_0$};
  \foreach \i/\lab in {0/NN,1/NT,2/TN,3/TT} {
    \draw[fill=ln-box] (1.4+1.25*\i,-0.3) rectangle ++(1.25,0.6);
    \node at (2.025+1.25*\i,0) {$C_{\mathrm{\lab}}$};
    \draw[->] (2.025+1.25*\i,-0.3) -- (2.025+1.25*\i,-1.1);
  }
  \node[font=\scriptsize, text=ln-muted, anchor=south] at (0.7,0.3) {\iflangja{履歴 2 ビット}{2-bit history}};
  \node[font=\scriptsize, text=ln-muted, anchor=south] at (3.9,0.3) {\iflangja{2 ビットカウンタ $\times 4$}{four 2-bit counters}};
  % 4-to-1 mux
  \draw[fill=white] (1.6,-1.1) -- (6.2,-1.1) -- (4.8,-1.7) -- (3.0,-1.7) -- cycle;
  \draw[->] (0.7,-0.3) -- (0.7,-1.4) -- (2.3,-1.4);
  \node[font=\scriptsize, below] at (1.3,-1.4) {\iflangja{選択}{select}};
  \draw[->] (3.9,-1.7) -- (3.9,-2.2) node[below, font=\scriptsize]
    {\iflangja{予測 = 選ばれたカウンタの上位ビット}{prediction = MSB of the selected counter}};
\end{tikzpicture}

  \caption{2 ビット履歴予測器のエントリ．履歴ビットが 4つの 2 ビット
    カウンタから 1つを選ぶ．分岐が解決されると，選ばれたカウンタが更新
    され，結果が履歴にシフトインされる．}
  \label{fig:l04-history}
\end{figure}

\begin{table}[htbp]
  \centering\small
  \caption{\cref{ex:l04-1bp-trace} のループ分岐に対する 2 ビット履歴
    予測器．履歴 NN，全カウンタ 00 から始める．履歴は古い結果から順に
    書き，予測ミスを $\times$ で示す．}
  \label{tab:l04-history-trace}
  \setlength{\tabcolsep}{3.6pt}
  \begin{tabular}{@{}lcccccccccccc@{}}
    \toprule
    実行           & 1 & 2 & 3 & 4 & 5 & 6 & 7 & 8 & 9 & 10 & 11 & 12 \\
    \midrule
    履歴           & NN & NT & TT & TN & NT & TT & TN & NT & TT & TN & NT & TT \\
    使うカウンタ   & 00 & 00 & 00 & 00 & 01 & 00 & 01 & 10 & 00 & 10 & 11 & 00 \\
    予測           & N & N & N & N & N & N & N & T & N & T & T & N \\
    結果           & T & T & N & T & T & N & T & T & N & T & T & N \\
    予測ミス       & $\times$ & $\times$ & & $\times$ & $\times$ & & $\times$ & & & & & \\
    \bottomrule
  \end{tabular}
\end{table}

\begin{definition}[$N$ ビット履歴予測器]
  \term[N びっとりれきよそくき]{$N$ ビット履歴予測器}[N-bit history predictor]%
  は，各分岐の直近 $N$ 回の結果と，エントリ当たり $2^N$ 個の 2 ビット
  カウンタ（可能な履歴ごとに 1つ）を保持する．1 エントリのコストは次の
  とおりである．
  \begin{equation}
    N + 2 \cdot 2^N \ \text{ビット}
    \label{eq:l04-history-cost}
  \end{equation}
\end{definition}

学習後，$N$ ビット履歴予測器は，周期が $N+1$ 以下のあらゆる繰り返し
パターンを正しく予測できる．そのようなパターンでは，直近 $N$ 回の結果が次の結果を
決めるからである．コストは $N$ に対して指数的に増大し，その大半は無駄に
なる．周期 $p$ の分岐は，$2^N$ 個のカウンタのうち高々 $p$ 個しか使わない．
また，使われるカウンタはそれぞれ個別に学習しなければならないので，履歴が
長いほど学習にも時間がかかる．

\begin{example}
  \label{ex:l04-history-cost}
  常に 10 回反復するループの分岐は，周期~10 のパターン
  T\,T\,T\,T\,T\,T\,T\,T\,T\,N を持つ．これを完全に予測するには $N = 9$，
  すなわち BHT エントリ当たり $9 + 2\cdot 512 = 1033$ ビットが必要であり，
  分岐が使うのは 512 個のカウンタのうち 10 個だけである．そのようなエントリ
  1024 個からなる BHT は 1 メガビットを超える．
\end{example}

\begin{keypoint}
  履歴により，予測器は偏りだけでなくパターンを学習できる．しかし，分岐
  ごとに専用の $2^N$ 個のカウンタを与えると，長い履歴は法外に高価になり，
  しかもそれらのカウンタのほとんどは一度も使われない．
\end{keypoint}

\section{共有カウンタ: PShare と GShare}
\label{sec:l04-share}

カウンタを分岐専用にしなければ，この無駄はなくなる．
\source{Lesson 4，動画 34--39；H\&P §3.3；\cite{mcfarling1993}．}BHT は
エントリごとに $N$ ビットの履歴だけを保持し，すべての分岐が 2 ビット
カウンタからなる 1つの大きなテーブルを共有する．

\begin{definition}[パターン履歴テーブル]
  \term[ぱたーんりれきてーぶる]{パターン履歴テーブル}[pattern history table]%
  \index{PHT|see{パターン履歴テーブル}}（PHT）とは，すべての分岐が共有する
  2 ビットカウンタのテーブルであり，分岐の履歴と PC の組合せ，典型的には
  それらのビットごとの XOR でインデックスされる．
\end{definition}

分岐が占めるカウンタは，その分岐が持つ異なる履歴の数だけになる．常に成立
する分岐の履歴は T\,T\,\ldots\,T であり，カウンタを 1つしか使わない．
交互に変わる分岐は 2つ，\cref{ex:l04-history-cost} の 10 回反復ループの
分岐は 10 個を使う．PC のビットを混ぜることで，同じ履歴を持つ異なる分岐が
異なるカウンタに送られる．衝突の可能性は残るが，PHT が大きければまれである．

\begin{example}
  12 ビットの履歴を保持する 1024 エントリの BHT は \num{12288} ビットを，
  $2^{12} = 4096$ 個のカウンタからなる PHT は \num{8192} ビットを占め，
  合計 \num{20480} ビットとなる．12 ビット履歴と専用カウンタを持つ
  1024 エントリの予測器では $1024 \times (12 + 2\cdot 4096) = \num{8400896}$
  ビットが必要であり，これは約 410 倍である．
\end{example}

\subsection{ローカル履歴とグローバル履歴}

PHT をインデックスする履歴の選び方は 2 通りある（\cref{fig:l04-pshare-gshare}）．

\begin{definition}[PShare]
  分岐の\term[ろーかるぶんきりれき]{ローカル分岐履歴}[private branch history]%
  とは，その分岐自身の直近の結果の系列であり，BHT エントリに保持される．
  \term[PShare よそくき]{PShare 予測器}[PShare predictor]は，ローカル履歴
  と PC の XOR で共有 PHT をインデックスする．
\end{definition}

\begin{definition}[GShare]
  \term[ぐろーばるぶんきりれき]{グローバル分岐履歴}[global branch history]%
  とは，すべての分岐の直近の結果の系列であり，解決された各分岐が結果を
  シフトインする単一のシフトレジスタに保持される．
  \term[GShare よそくき]{GShare 予測器}[GShare predictor]は，グローバル
  履歴と PC の XOR で共有 PHT をインデックスする．
\end{definition}

\begin{figure}[htbp]
  \centering
  % PShare (private history) and GShare (global history), both with a shared
% pattern history table of two-bit counters.
\begin{tikzpicture}[x=1cm,y=1cm, font=\footnotesize\sffamily, >={Stealth[length=1.8mm]},
  box/.style={draw, fill=ln-box, minimum height=0.6cm, minimum width=0.9cm, inner sep=2pt},
  xor/.style={inner sep=0pt, font=\large}]
  % ---------------- (a) PShare
  \node[anchor=west, font=\footnotesize\bfseries] at (-0.45,1.45) {(a) PShare};
  \node[box] (pca) at (0,0) {PC};
  \draw[fill=ln-box] (1.3,-0.7) rectangle (2.9,0.7);
  \foreach \y in {-0.35,0.35} { \draw[ln-rule] (1.3,\y) -- (2.9,\y); }
  \fill[ln-accent!15] (1.31,-0.34) rectangle (2.89,0.34);
  \node[font=\scriptsize, text=ln-muted] at (2.1,0.95) {\iflangja{BHT（ローカル履歴）}{BHT (private histories)}};
  \draw[->] (pca.east) -- (1.3,0);
  \node[xor] (xa) at (4.0,0) {$\oplus$};
  \draw[->] (2.9,0) -- (xa.west);
  \draw[->] (pca.south) -- (0,-1.0) -- (4.0,-1.0) -- (xa.south);
  \draw[fill=ln-box] (5.0,-0.7) rectangle (6.2,0.7);
  \foreach \y in {-0.35,0,0.35} { \draw[ln-rule] (5.0,\y) -- (6.2,\y); }
  \node[font=\scriptsize, text=ln-muted] at (5.6,0.95) {\iflangja{PHT（2 ビットカウンタ）}{PHT (2-bit counters)}};
  \draw[->] (xa.east) -- (5.0,0);
  \draw[->] (6.2,0) -- (7.0,0) node[right] {\iflangja{予測}{prediction}};
  % ---------------- (b) GShare
  \begin{scope}[yshift=-3.4cm]
  \node[anchor=west, font=\footnotesize\bfseries] at (-0.45,1.45) {(b) GShare};
  \foreach \i/\b in {0/1,1/0,2/1,3/1,4/0,5/1} {
    \draw[fill=ln-accent!15] (1.3+0.2667*\i,-0.3) rectangle ++(0.2667,0.6);
    \node[font=\scriptsize\ttfamily] at (1.4333+0.2667*\i,0) {\b};
  }
  \node[font=\scriptsize, text=ln-muted, align=center] at (2.1,0.8) {\iflangja{グローバル履歴\\（シフトレジスタ）}{global history\\(shift register)}};
  \node[xor] (xb) at (4.0,0) {$\oplus$};
  \draw[->] (2.9,0) -- (xb.west);
  \node[box] (pcb) at (0,-1.0) {PC};
  \draw[->] (pcb.east) -- (4.0,-1.0) -- (xb.south);
  \draw[fill=ln-box] (5.0,-0.7) rectangle (6.2,0.7);
  \foreach \y in {-0.35,0,0.35} { \draw[ln-rule] (5.0,\y) -- (6.2,\y); }
  \node[font=\scriptsize, text=ln-muted] at (5.6,0.95) {\iflangja{PHT（2 ビットカウンタ）}{PHT (2-bit counters)}};
  \draw[->] (xb.east) -- (5.0,0);
  \draw[->] (6.2,0) -- (7.0,0) node[right] {\iflangja{予測}{prediction}};
  \end{scope}
\end{tikzpicture}

  \caption{共有カウンタを用いる予測器．(a) PShare は分岐自身の履歴を PC と
    組み合わせ，(b) GShare は全分岐のグローバル履歴を PC と組み合わせる．}
  \label{fig:l04-pshare-gshare}
\end{figure}

PShare は，独自のパターンを持つ分岐，すなわち交互に変わる分岐や短いループ
に適している．GShare は，直前に実行された他の分岐の結果に依存する
\emph{相関のある}分岐に適している．

\needspace{9\baselineskip}
\begin{example}
  次の断片では，2 番目の分岐は，1 番目の分岐が不成立のときに成立し，
  成立のときに不成立となる．
  \begin{lstlisting}[language={[x86masm]Assembler}, numbers=none]
        BNE  R1, R0, Skip1   ; R1 == 0 でなければスキップ
        ADDI R2, R2, 1       ; ゼロの個数を数える
Skip1:  BEQ  R1, R0, Skip2   ; R1 == 0 ならスキップ
        ADD  R3, R3, R1      ; 非ゼロの値を合計
Skip2:  ...
\end{lstlisting}
  R1 の値がデータに依存するなら，どちらの分岐のローカル履歴もランダムに
  見え，PShare はうまく予測できない．\margin{結果がデータに依存する分岐は，
  予測するより取り除く方が安上がりなことがある．第5章はそのような分岐を
  プレディケーションで置き換える．}2 番目の分岐から見たグローバル履歴には
  1 番目の分岐の結果が含まれるので，これらの履歴に対応する PHT カウンタの
  学習が済めば，PHT でのエイリアシングがない限り，GShare は 2 番目の分岐を
  完全に予測する．
\end{example}

どちらもすべての分岐に対して優れているわけではない．\caution{標準的な名称は
GShare（McFarling）だけである．講義が PShare と呼ぶものは，文献ではローカル
履歴予測器（Yeh と Patt の PAs 方式）と呼ばれる．}GShare も原理的には
分岐自身のパターンを学習できるが，その間に実行された他の分岐の結果が
グローバル履歴に混ざり込むので，同じパターンに PShare よりかなり長い履歴
---そしてより大きな PHT---が必要になる．プログラムには両方の種類の分岐が
含まれるので，自然な次の一歩は両方の予測器を使うことである．

\section{トーナメント予測器と階層型予測器}
\label{sec:l04-hybrid}

分岐によって最適な予測器が異なるので，プロセッサは複数の予測器を組み合わせ，
分岐ごとにそれに適したものを使うことができる．\source{Lesson 4，動画 40--43；
H\&P §3.3；\cite{mcfarling1993}．}組合せ方には対称なものと階層的なものが
ある．

\subsection{トーナメント予測器}

\term[とーなめんとよそくき]{トーナメント予測器}[tournament predictor]は，
2つの予測器，例えば GShare と PShare を並行して動作させ，どちらを信頼
するかを分岐ごとに個別に決める．

\begin{definition}[メタ予測器]
  トーナメント予測器の\term[めたよそくき]{メタ予測器}[meta-predictor]とは，
  PC でインデックスされる 2 ビットカウンタのテーブルであり，ある分岐に
  対応するカウンタが，2つの構成予測器のどちらが正しいかを予測する
  （\cref{fig:l04-tournament}）．
\end{definition}

両方の構成予測器は，通常どおりすべての分岐で学習する．\margin{Alpha 21264 は，10 ビットのローカル履歴 1024 個，12 ビットのグローバル
履歴で引く 4096 エントリの表，4096 エントリの選択器を組み合わせ，合計約
29 キロビットであった（H\&P §3.3）．}メタ予測器のカウンタ
が変化するのは構成予測器の正誤が分かれたときだけであり，そのときは正しかった
方へ向かって動く．両方とも正しいとき，または両方とも誤ったときは変化しない．

\begin{figure}[htbp]
  \centering
  % Tournament predictor: a meta-predictor chooses between two predictors.
\begin{tikzpicture}[x=1cm,y=1cm, font=\footnotesize\sffamily, >={Stealth[length=1.8mm]},
  box/.style={draw, fill=ln-box, minimum height=0.7cm, minimum width=2.5cm, inner sep=2pt, align=center}]
  \node[draw, fill=ln-box, minimum width=0.9cm, minimum height=0.6cm] (pc) at (0,0) {PC};
  \node[box] (p1) at (3.0,1.2) {PShare};
  \node[box] (p2) at (3.0,0) {GShare};
  \node[box, fill=ln-accent!15] (meta) at (3.0,-1.3) {\iflangja{メタ予測器}{meta-predictor}\\[-1pt]{\scriptsize\iflangja{（2 ビットカウンタ）}{(2-bit counters)}}};
  \draw (pc.east) -- (1.2,0);
  \draw (1.2,1.2) -- (1.2,-1.3);
  \fill (1.2,0) circle (1.2pt);
  \draw[->] (1.2,1.2) -- (p1.west);
  \draw[->] (1.2,0) -- (p2.west);
  \draw[->] (1.2,-1.3) -- (meta.west);
  % mux
  \draw[fill=white] (5.4,1.6) -- (6.0,1.2) -- (6.0,-0.4) -- (5.4,-0.8) -- cycle;
  \draw[->] (p1.east) -- (5.4,1.2);
  \draw[->] (p2.east) -- (5.4,0);
  \draw[->] (meta.east) -- (5.7,-1.3) -- (5.7,-0.6);
  \node[font=\scriptsize, below] at (5.0,-1.3) {\iflangja{どちら?}{which one?}};
  \draw[->] (6.0,0.4) -- (6.8,0.4) node[right] {\iflangja{予測}{prediction}};
\end{tikzpicture}

  \caption{トーナメント予測器．メタ予測器が，どの構成予測器の予測を使うか
    を分岐ごとに選ぶ．}
  \label{fig:l04-tournament}
\end{figure}

\subsection{階層型予測器}

トーナメント予測器は，よい---したがって高価な---予測器を 2つ組み合わせ，
すべての分岐で両方を更新する．しかし，ほとんどの分岐は予測が容易である．

\begin{definition}[階層型予測器]
  \term[かいそうがたよそくき]{階層型予測器}[hierarchical predictor]は，
  すべての分岐を扱う安価な予測器と，安価な予測器がうまく扱えない分岐に
  対してだけ使用・学習される，より優れた高価な 1つ以上の予測器とを
  組み合わせる．
\end{definition}

高価な予測器が学習すべき分岐ははるかに少ないので，同じ精度ならトーナメント
予測器の場合よりずっと小さくでき，同じサイズなら精度をより高くできる．
\margin{AMD の Zen~2 は速度で階層化する．高速なパーセプトロンが予測し，低速な
TAGE 予測器がそれを上書きする．予測ミスは Zen~1 より約 \SI{30}{\percent}
少ない（AMD）．}

典型的な構成は 3 レベルからなる（\cref{fig:l04-hierarchical}）．PC で
インデックスされる 2 ビットカウンタの大きなテーブルが，すべての分岐を予測
する．その上に，小さなローカル履歴予測器と小さなグローバル履歴予測器が
あり，選ばれた分岐のエントリだけを保持する．それらのエントリには PC の
ビットからなる\emph{タグ}が付けられており，参照時にその分岐のエントリが
あるかどうかが分かる．予測は，一致するエントリを持つ最も上位のレベルから
得る．カウンタテーブルは常に更新され，分岐に一致するローカルまたは
グローバルのエントリも同様に更新される．カウンタテーブルが予測ミスした
分岐にはローカル予測器のエントリが与えられ，ローカル予測器が予測ミスした
分岐にはグローバル予測器のエントリが与えられる．カウンタでよく予測できる
分岐が高価なレベルのエントリを占めることはまれなので，それらの小さな
テーブルは，エントリを必要とする分岐のために残される．\sidenote{Intel の
Pentium~M（2003）が典型例である．全分岐を対象とする 2 ビットカウンタの表を，
タグ付きのローカル履歴予測器とグローバル履歴予測器，ループ分岐の反復回数を
学習するループ検出器，そして間接分岐専用の予測器が補う（Gochman ら，
\emph{Intel Technology Journal}~7(2)，2003）．}

\begin{figure}[htbp]
  \centering
  % Hierarchical predictor: tagged upper levels override a cheap 2BC table.
\begin{tikzpicture}[x=1cm,y=1cm, font=\footnotesize\sffamily, >={Stealth[length=1.8mm]},
  box/.style={draw, fill=ln-box, minimum height=0.8cm, minimum width=3.0cm, inner sep=2pt, align=center}]
  \node[draw, fill=ln-box, minimum width=0.9cm, minimum height=0.6cm] (pc) at (0,0) {PC};
  \node[box, fill=ln-accent!15] (g) at (3.2,1.3) {\iflangja{グローバル履歴予測器}{global-history predictor}\\[-1pt]{\scriptsize\iflangja{小・タグ付き}{small, tagged}}};
  \node[box, fill=ln-accent!8] (l) at (3.2,0) {\iflangja{ローカル履歴予測器}{local-history predictor}\\[-1pt]{\scriptsize\iflangja{小・タグ付き}{small, tagged}}};
  \node[box] (b) at (3.2,-1.3) {\iflangja{2 ビットカウンタテーブル}{2-bit counter table}\\[-1pt]{\scriptsize\iflangja{大・タグなし}{large, untagged}}};
  \draw (pc.east) -- (1.2,0);
  \draw (1.2,1.3) -- (1.2,-1.3);
  \fill (1.2,0) circle (1.2pt);
  \foreach \n in {g,l,b} { \draw[->] (1.2,0 |- \n.west) -- (\n.west); }
  \node[draw, fill=white, align=left, font=\scriptsize, minimum height=3.4cm, text width=2.3cm] (sel) at (6.8,0)
    {\iflangja{1.\ グローバルがヒット\\\quad$\to$ それを使う\\[3pt]2.\ ローカルがヒット\\\quad$\to$ それを使う\\[3pt]3.\ それ以外\\\quad$\to$ 2BC を使う}{1.\ global hit\\\quad$\to$ use global\\[3pt]2.\ local hit\\\quad$\to$ use local\\[3pt]3.\ otherwise\\\quad$\to$ use 2BC}};
  \foreach \n in {g,l,b} { \draw[->] (\n.east) -- (sel.west |- \n.east); }
  \draw[->] (sel.east) -- ++(0.6,0) node[right] {\iflangja{予測}{prediction}};
\end{tikzpicture}

  \caption{階層型予測器．タグ付きの上位レベルは，保持している分岐について
    2 ビットカウンタテーブルの予測を置き換える．}
  \label{fig:l04-hierarchical}
\end{figure}

\cref{tab:l04-which} は，どの種類の分岐をどの種類の予測器が扱うかをまとめた
ものである．

\begin{table}[htbp]
  \centering\small
  \caption{分岐の振る舞いと，それをうまく扱える最も単純な予測器．}
  \label{tab:l04-which}
  \begin{tabular}{@{}p{0.52\linewidth}p{0.4\linewidth}@{}}
    \toprule
    分岐の振る舞い & 予測器 \\
    \midrule
    ほとんど常に同じ方向 & 2 ビットカウンタ \\
    独自の繰り返しパターン（交互，短いループ） & ローカル履歴（PShare） \\
    先行する分岐と相関がある & グローバル履歴（GShare） \\
    1つのプログラム内で上記が混在 & トーナメントまたは階層型 \\
    関数からのリターン & リターンアドレススタック（次節） \\
    \bottomrule
  \end{tabular}
\end{table}

\section{リターンアドレススタック}
\label{sec:l04-ras}

関数からのリターン（\texttt{RET}）は無条件分岐なので，方向の予測は自明で
あり，常に成立と予測すればよい．\source{Lesson 4，動画 44--47；H\&P §3.9；
\cite{kaeli1991}．}しかし，その分岐先は関数がどこから呼ばれたかに依存する．
BTB は最新の分岐先しか覚えていないので，関数が異なる場所から代わる代わる
呼ばれると，リターンは毎回予測ミスする．

\begin{example}
  ある関数が 2 か所から呼ばれる．
  \begin{lstlisting}[language={[x86masm]Assembler}, numbers=none]
0x1000         CALL Func    ; 戻りアドレス 0x1004
  ...
0x1040         CALL Func    ; 戻りアドレス 0x1044
  ...
0x2000  Func:  ...
0x2030         RET          ; まず 0x1004 へ，次に 0x1044 へ
\end{lstlisting}
  1 回目のリターンが解決されると，\texttt{RET} の BTB エントリには
  \texttt{0x1004} が入る．このリターン自体は，エントリがまだ空だったので
  予測ミスする．\texttt{0x1044} へ戻る 2 回目のリターンは予測ミスし，
  エントリを上書きする．2 か所からの呼び出しが交互に起こると，リターンは
  毎回予測ミスする．
\end{example}

\begin{definition}[リターンアドレススタック]
  \term[りたーんあどれすすたっく]{リターンアドレススタック}[return address stack]%
  \index{RAS|see{リターンアドレススタック}}（RAS）とは，フェッチステージに
  ある小さなハードウェアスタックである．コールがフェッチされると，その
  戻りアドレス（PC+4）をプッシュする．リターンがフェッチされると，先頭の
  エントリをポップし，予測した分岐先として用いる．
\end{definition}

コールとリターンは入れ子になるので，RAS はリターンをほぼ完全に予測する
（\cref{fig:l04-ras}）．RAS は小さいので，コールが深く入れ子になると満杯に
なりうる．\margin{実際のリターンアドレススタックは小さい．Intel のコアでは
ハードウェアスレッド当たり 16 エントリ，AMD の Zen~2 では 32 エントリで
ある．}RAS が満杯のときにコールがフェッチされると，プッシュはラップ
アラウンドして最も古いエントリを上書きする．リターンはコールの逆順に起こるので，最も早く
必要になるのは最新のエントリであり，またほとんどのコールとリターンは
コールツリーの深い位置で起こる．最も古いエントリを上書きしても，失われる
のは最も外側の関数からのリターンの予測だけであり，それはまれである．一方，
新しい戻りアドレスを捨てれば，最も深いレベルで頻繁に起こるコールの予測が
失われる．

\begin{figure}[htbp]
  \centering
  % Return address stack: calls push, returns pop; a full stack wraps around.
\begin{tikzpicture}[x=1cm,y=1cm, font=\footnotesize\sffamily, >={Stealth[length=1.8mm]}]
  \foreach \i/\v in {0/0x1004,1/0x2210,2/0x2A08,3/0x3104} {
    \draw[fill=ln-box] (0,0.55*\i) rectangle (1.8,0.55*\i+0.55);
    \node[font=\footnotesize\ttfamily] at (0.9,0.55*\i+0.275) {\v};
  }
  \draw[very thick, ln-accent] (0,1.65) rectangle (1.8,2.2);
  \node[font=\scriptsize, text=ln-muted, anchor=north] at (0.9,-0.05) {\iflangja{最も古い}{oldest}};
  % pop from the top
  \draw[->, thick] (1.85,1.925) -- (3.0,1.925);
  \node[anchor=west, font=\scriptsize] at (3.05,1.925) {\iflangja{RET: 先頭をポップ → 予測した分岐先}{RET: pop top $\to$ predicted target}};
  % push onto a full stack wraps around to the oldest entry
  \draw[thick] (3.0,2.6) -- (0.9,2.6);
  \draw[->, thick, dashed] (0.9,2.6) -- (-0.4,2.6) to[out=180,in=180,looseness=1.6] (-0.05,0.275);
  \node[anchor=west, font=\scriptsize] at (3.05,2.6) {\iflangja{CALL: 戻りアドレス（PC+4）をプッシュ}{CALL: push return address (PC+4)}};
  \node[anchor=east, align=right, font=\scriptsize, text=ln-muted] at (-1.45,1.35)
    {\iflangja{満杯のとき:\\プッシュで\\最古を上書き}{when full:\\push overwrites\\oldest}};
\end{tikzpicture}

  \caption{4 エントリのリターンアドレススタック．コールでプッシュし，
    リターンでポップする．スタックが満杯のときは，プッシュがラップアラウンド
    して最も古いエントリを上書きする．}
  \label{fig:l04-ras}
\end{figure}

\subsection{リターンの識別}

RAS は，\texttt{RET} がデコードされる前，すなわちフェッチされている最中に
使わなければならない．命令がリターンであることを知る方法は 2つある．
\begin{itemize}
  \item \textbf{予測する．}予測器のテーブルに，PC ごとにその命令がリターン
    かどうかを格納する．エントリは，その命令が初めてデコードされたときに
    設定される．あるアドレスにある命令は変化しないので，この予測はほぼ
    常に正しい．
  % この項目を章のまとめと同じページに置く（まとめだけのページを避ける）
  \needspace{9\baselineskip}
  \item \textbf{プリデコードする．}\term[ぷりでこーど]{プリデコード}[predecoding]%
    では，プロセッサは各命令をメインメモリからプロセッサのキャッシュに
    取り込む際に部分的にデコードし，例えばそれが分岐かリターンかといった
    数ビットの追加情報を，キャッシュされた命令とともに格納する．フェッチ
    ステージはこれらのビットを命令と一緒に読むので，命令の種類がただちに
    分かる．
\end{itemize}

\begin{keypoint}[章のまとめ]
  すべてのフェッチは予測である．予測ミスは，それがフラッシュするすべての
  命令の作業を無駄にし（\cref{eq:l04-relative}），フラッシュされる命令数は
  パイプラインの深さと幅とともに増える．不成立予測はコストがかからないが
  不正確である．BTB は分岐先を，BHT は方向を予測し，BHT は 2 ビット飽和
  カウンタを用いると最もよく機能する．履歴ベース予測器はパターンを学習する．
  そのカウンタを PHT で共有すれば長い履歴も現実的なコストで実現でき，自己反復する
  分岐にはローカル履歴（PShare），相関のある分岐にはグローバル履歴
  （GShare）を用いる．トーナメント予測器と階層型予測器は複数の予測器を
  組み合わせ，リターンアドレススタックは関数からのリターンを予測する．
\end{keypoint}

\end{document}
